(1) 近轴近似下两缝至屏上 $x$ 的光程差 $\Delta\ell\simeq dx/L$，相位差为 $2\pi dx/(\lambda L)$，因此明纹相邻间距 $\Delta x=\lambda L/d$。此式只取决于双缝间距、波长和观测几何，暂不涉及照明的空间相干长度。
(2) 将线光源坐标记为 $s\in[-a/2,a/2]$。不同 $s$ 发出的光彼此非相干，而每一光源点到双缝的相位差为 $2\pi ds/(\lambda S)$；归一化相干度为 $\gamma_{12}=(1/a)\int_{-a/2}^{a/2}e^{i2\pi ds/(\lambda S)}ds=\sin(\pi ad/(\lambda S))/[\pi ad/(\lambda S)]$。
(3) 若各缝在观察屏贡献强度为 $I_s$，则 $I(x)=2I_s[1+\gamma_{12}\cos(2\pi dx/(\lambda L))]$，此处 $\gamma_{12}$ 为实数，若为负则对应 $\pi$ 相位翻转。可见度 $\mathcal V=|\gamma_{12}|$，其第一个非平凡零点为 $ad/(\lambda S)=1$，即 $a_{\rm zero}=\lambda S/d$。
(4) 当 $a\to0$ 时利用 $\sin z/z\to1$，恢复点源的完全空间相干照明；$d\to0$ 时两采样位置重合，相干度也趋一，但条纹间距发散而不再形成可辨周期条纹。线光源中不同发光点给出的条纹整体横向错位，强度非相干叠加使对比度被平均掉，并非每个点源自身失去干涉能力。
